\section{Введение}
\label{sec:Chapter0} \index{Chapter0}

%В этой части надо описать предметную область, задачу из которой вы будете
%решать, объяснить её актуальность (почему надо что-то делать сейчас?). Здесь же
%стоит ввести определения понятий, которые вам понадобятся в постановке задачи.

Большая часть современных криптографических протоколов базируются на сложности задач факторизации натурального числа и дискретного логарифмирования элемента в какой-то абелевой группе.

Такими являются, например, широко известный и много применявшийся раньше $\mathsf{RSA}$ ({\it Rivest–Shamir–Adleman cryptosystem}), а также современные российские ГОСТы ({\it государственные стандарты}) для цифровой подписи документов и получения общего секретного ключа.

При этом в $1994$\;-\,$1997$ годах Шор разработал алгоритм на квантовом компьютере, который решает обе эти задачи в $\N$ за полиномиальное время от длины числа~\cite{Shor_1997}. В дальнейшем алгоритм улучшался и модифицировался, в том числе, чтобы работать на эллиптических кривых~\cite{10.5555/2011528.2011531}.

Сейчас многими компаниями и государствами ведутся исследования по построению мощных стабильных квантовых компьютеров. В случае успеха оных все вышеупомянутые протоколы окажутся уязвимыми.

В связи с этим, начиная с $2016$\;-\:$2017$ годов NIST ({\it National Institute of Standards and Technology в США}) проводит конкурс по отбору лучших протоколов постквантовой ({\it устойчивой к атакам в том числе на квантовом компьютере}) криптографии.

Немалая доля из них основана на решёточной криптографии вообще и $\mathsf{LWE}$-({\it Learning With Errors})-problem в частности~\cite{10.1145/1060590.1060603}. Одним из главных преимуществ здесь является $\mathsf{ACH}$ ({\it Average-Case}), а не $\mathsf{WCH}$ ({\it Worst-Case}) сложность ({\it Hardness}). Это значит, что трудными для взлома будут случайные, а не худшие из генерируемых параметры протокола. 

Есть несколько широко известных $\mathsf{WCH}$ проблем на решётках: $\mathsf{SVP}$ ({\it Shortest Vector Problem}), $\mathsf{GapSVP}$, $\mathsf{SIVP}$ ({\it Shortest independent vectors problem}) и т.п. В $1996$ году Ажтай свёл часть из них к $\mathsf{ACH}$ проблеме $\mathsf{SIS}$ ({\it Short Integer Solution})~\cite{10.1145/237814.237838}. Позже Регев проделал аналогичную вещь с $\mathsf{LWE}$.

К сожалению, протоколы, основанные на этой задаче, хотя являются простыми и понятными, требуют довольно много вычислений. Поэтому с целью оптимизации в статье~\cite{10.1145/2535925} вводится $\mathsf{RLWE}$-({\it Ring LWE})-problem. Это модификация $\mathsf{LWE}$ в многочленах над $\Z_q$ c ускоренными вычислениями за счёт $\mathsf{NTT}$ ({\it Number Theoretic Transform}).

Более подробно все эти и некоторые дальнейшие концепции описаны в вышедшем в $2022$ году обзоре~\cite{Li2022ATI}. В моей же работе я буду использовать и улучшать результаты следующей статьи~\cite{9457596}.

В ней процесс $\mathsf{RLWE}$-шифрования рассматривается, как передача данных по зашумлённому $\mathsf{RLWE}$-каналу. Исследуются свойства этого канала, описывается, как увеличить скорость передачи и при этом понизить $\mathsf{DFR}$ ({\it Decryption Failure Probability/Rate}) при расшифровке сообщения.

Для уменьшения $\mathsf{DFR}$ при фиксированном $\mathsf{CFR}$ ({\it Coefficient Failure Rate}) можно использовать $\mathsf{ECC}$ ({\it Error Correcting Codes}), исправляющие $\leq\!\?t$ символьных неточностей при передаче сообщения.

В моей работе я пробую разные способы мерить кол-во ошибок (Хэмминг, Ли) и вычисляю, как изменяется от этого скорость $\mathsf{RLWE}$-канала.
